\clearpage
\fancyhead[L]{\small SOLDNER. AZIMUTHE ZU REDUCIREN.}
\SourceLeaf{592}
\begin{center}\small Göttingische gelehrte Anzeigen. Issue 46. Pages 449..456. 23 March 1815.\\[3pt]\rule{.72\linewidth}{.3pt}\end{center}

I. \emph{A New Method for Reducing Observed Azimuths. By} J. \textsc{Soldner}. Munich. 14 pages in quarto. (Printed separately from the memoirs of the Munich Academy for 1813.)

II. \emph{Contributions to the Calculation of Observed Azimuths by} \textsc{Anton von Stefenelli}. Same place. Printed with Storno's type. 30 pages in quarto.

Because azimuths change at an uneven rate, we cannot equate the mean of a series of azimuth differences between a terrestrial object and a celestial body with the azimuth difference corresponding to the mean of all the observation times. Of course, if a reading is taken after each individual observation, one will not shrink from calculating each one separately to obtain as many separate results. As a rule, though, one takes a reading only at the end of a series so that one can make more observations instead. Calculating a large number of observations one by one would be a forbiddingly lengthy task. Confusing the mean of all the observed azimuth differences with the azimuth difference at the mean of all the observation times would produce errors that cannot be overlooked given the precision of observations made with theodolites such as those now supplied by \textsc{Reichenbach}. Partly for this reason, a famous practical astronomer proposed making these observations only around the time of culmination. The change in azimuth can indeed be treated as uniform then, and one has the further advantage that errors in the elevation of the pole or in the declination of the celestial body have an altogether imperceptible effect on the result. But this method has other drawbacks in turn. The greater the altitude of the body observed, the greater the effect of small errors in adjusting the instrument. The adjustment itself is already very delicate and easily disturbed again to some degree if the instrument cannot be set up very firmly. An error in determining the time affects the azimuth more adversely here than at low altitudes. In observing the Sun, one is limited to the winter months at our latitudes or farther south, insofar as the telescope can be tilted only up to $40^{\circ}$. It is true that fixed stars of the first magnitude can be observed even in daylight with those instruments. But if the time must be determined from corresponding solar altitudes, any error in the solar tables must also be taken into account. If, on the other hand, one has a good transit instrument, azimuth observations will in most cases be entirely unnecessary, since that instrument makes it easy to place an artificial object in the meridian plane with great precision. For these reasons, we would give preference to azimuth determinations away from culmination and at moderate altitudes. If there is still a slight uncertainty in the elevation of the pole or in the declination of the celestial body, one can readily make it harmless simply by combining observations made on both sides of the meridian and at nearly equal distances from it. Nor need the uneven change in azimuth mentioned earlier matter here, since it can easily be taken into account with the greatest precision. Indeed, it is clear from \textsc{Taylor}'s theorem that the azimuth corresponding to hour angle $t+\theta$ can be expanded in a series $A+B\theta+C\theta\theta+D\theta^3+\text{etc.}$, of which a few terms will suffice as long as $\theta$ stays within narrow bounds. If a total of $n$ observations have been made at the hour angles $t+\theta$, $t+\theta'$, $t+\theta''$, $t+\theta'''$, and so on, then the mean of all the azimuths is \SourceLeaf{593}clearly
\[
\begin{aligned}&= A + \frac{\theta+\theta'+\theta''+\theta'''+\text{etc.}}{n}\cdot B\\[3pt]
&+ \frac{\theta\theta+\theta'\theta'+\theta''\theta''+\theta'''\theta'''+\text{etc.}}{n}\cdot C\\[3pt]
&+ \text{etc.}\end{aligned}
\]
Since the choice of a base value of $t$ is arbitrary here, it is simplest to choose the value corresponding to the mean of all the observation times. This gives $\theta+\theta'+\theta''+\theta'''+\text{etc.}=0$, so the mean of all the azimuths contains the following form.
\[
A + \frac{1}{n}\,C\,\Sigma\theta\theta + \frac{1}{n}\,D\,\Sigma\theta^3
\]
In practice there is never any need to go beyond the third-order terms. Indeed, in the great majority of cases the second-order correction $\frac{1}{n}\,C\,\Sigma\theta\theta$ will suffice, and one can clearly use
\[
2\,C \times \frac{\Sigma\,2\sin\frac{1}{2}\theta^2}{n}
\]
in its place, since the error is only of the fourth order. The advantage is that one can use the familiar auxiliary table for reducing circummeridian altitudes, which gives $2\sin\frac{1}{2}\theta^2$ directly in seconds. It is clear that in the third term, with $D$ as it comes from the expansion by \textsc{Taylor}'s theorem, one must also include the factor $\frac{15^3}{206265^2}$ or $\frac{900^3}{206265^2}$ or $\frac{9000^3}{206265^2}$ when $\theta$ is expressed in seconds of time, minutes of time, or tens of minutes of time, respectively.

This briefly described procedure is the one Mr. \textsc{Soldner} presents in the paper under review. It is so simple and suggests itself so readily that one must wonder how several practical astronomers missed it in the same or quite similar circumstances. That is precisely why Mr. \textsc{Soldner} deserves the thanks of practical astronomers for making it known. All the more so because the expansions are carried out skillfully and elegantly, and the final results, that is, the expressions for the coefficients $B$, $C$, $D$ (for Mr. \textsc{Soldner} took the expansion this far), are put into a very convenient form.

\begingroup\rightskip=0pt plus 3em
The reviewer has used this very procedure for frequent azimuth observations made over the past couple of years with a \textsc{Reichenbach} theodolite, though without taking into account the third-order term, which is usually superfluous in practice. He hopes that astronomical readers will be glad to find a few more remarks on the subject here. He used the formula
\[
2\,C = -\,\frac{\cos\varphi\cdot \operatorname{tang}\,\delta\cdot \sin A^2}{\sin t}
\left(1 + \frac{2\cos\varphi\cdot \sin 2A}{\sin t\cdot \sin 2\delta}\right)
\]
where $\varphi$, $\delta$, $t$, $A$ denote the elevation of the pole, declination, hour angle and azimuth. The formula can also be written as follows.
\[
a = \frac{\cos\varphi\cdot \operatorname{tang}\,\delta\cdot \sin A^2}{\sin t}, \qquad
2\,C = -\,a\left(1 + \frac{2a}{\operatorname{tang}\,A\cdot \sin\delta^2}\right)
\]
This formula is in fact identical to one also given by Mr. \textsc{Soldner} on p. 12, which he likewise recommends only when one intends to leave the third-order term out of the calculation. The paper leaves unmentioned that \emph{under this assumption} it is more convenient to calculate the azimuth itself by the following formulas, which are familiar in their own right, in place of \textsc{Neper}'s.
\SourceLeaf{594}
\[
\begin{gathered}
\frac{\operatorname{tang}\delta}{\cos t}=\operatorname{tang}\psi\\[5pt]
\frac{\cos\psi\cdot\operatorname{tang}t}{\sin(\varphi-\psi)}=\operatorname{tang}A
\end{gathered}
\]
If the third-order term is to be taken into account, however, for $C$ and $D$ one will use Mr. \textsc{Soldner}'s formulas on p. 9, and a small modification to their form can also be made, to be mentioned below.\par\endgroup

Whenever the second-order correction suffices, even the calculation of the coefficient $C$ can be avoided altogether by a small trick. It can be used to advantage on various other occasions as well, and we will describe it here. As before, let $t$ be the hour angle at the mean of all the observation times, or $\theta+\theta'+\theta''+\theta'''+\text{etc.}=0$, and let
\[
\tau=\BareRoot{\frac{\theta\theta+\theta'\theta'+\theta''\theta''+\theta'''\theta'''+\text{etc.}}{n}}
\]
The mean of all $n$ azimuths is then to be set equal to the mean of the two azimuths corresponding to the hour angles $t-\tau$ and $t+\tau$. The two azimuths can be calculated directly about as easily as the single azimuth $A$ and the coefficient $C$. To find $\tau$ more easily, one may use the reduction table already mentioned without hesitation. The entry in that table for $\tau$ expressed in time will be the mean of the individual $n$ quantities corresponding in the table to the arguments $\theta$, $\theta'$, $\theta''$, $\theta'''$, and so on, expressed in time. The procedure also has the advantage that it lets one take various incidental circumstances into account without any further work, such as the change in the Sun's declination in the present problem when the Sun is the object in question. The trick, however, applies directly only if the hour angles are known and the azimuths are sought. When observed azimuth differences are used to determine the time, additional aids would have to be brought in. This is not the place to explain them.

Viewed from a purely mathematical standpoint, the investigation Mr. \textsc{Soldner} undertakes in fact solves the following problem. If two sides of a spherical triangle are fixed, express the finite changes of an adjacent angle as a series in successive powers of the finite changes of the included angle. The counterpart would be to represent the finite changes of the third side in a similar way under the same conditions. Its solution is likewise necessary on several occasions for the precise reduction of observations with multiplying instruments. Paper No. II deals with these two problems and various applications. Its author was already familiar with Mr. \textsc{Soldner}'s paper, as one can see from various remarks about it, some of them quite inappropriate. One is the strange assertion on p. 29 that \textsc{Soldner}'s procedure is nothing more than the method used by \textsc{Svanberg}. Anyone who compares the two methods will see so clearly how groundless that judgment is that we need spend no longer on it. As for the execution, we cannot give \textsc{Stefenelli}'s paper the same praise as \textsc{Soldner}'s. Both series by which Mr. \textsc{Stefenelli} thinks he solves the two problems are \emph{incorrect} already in the third-order terms. If we remember rightly, the two series were announced some time ago in a scholarly journal under the inappropriate heading of new discoveries. The fact that these terms are usually insignificant in practice does not excuse the author. He was free to stop at the second-order terms, but once he undertook to go \SourceLeaf{595}further, one expects correct results. All the more so because Mr. \textsc{Soldner} had already fully solved one problem through the third order, and because the entire merit of these expansions, which are in themselves quite easy, lies solely in their correctness and ease of use. The only thing we can say in praise of \textsc{Stefenelli}'s paper is that the still-correct second-order term (in the first problem, which is Mr. \textsc{Stefenelli}'s second) has a somewhat more convenient form than in Mr. \textsc{Soldner}'s paper. But Mr. \textsc{Stefenelli} must have failed to notice that \textsc{Soldner}'s series can at once be put into the same form by the following simple substitutions based on familiar equations.
\[
\begin{gathered}
\sin\varphi+\sin\delta=\frac{2\cos\frac{1}{2}z^2\cdot\sin2\gamma}{\sin t}\\[5pt]
\sin\varphi-\sin\delta=\frac{2\sin\frac{1}{2}z^2\cdot\sin2\beta}{\sin t}
\end{gathered}
\]
Had he noticed this, comparing the third-order term could have shown him that his own was wrong. It is hard to see how he could have overlooked that confusing arcs with their sines and tangents on p.\,3 was bound to introduce errors already at the third order.

\begingroup\rightskip=0pt plus 3em
According to a closing remark, Mr. \textsc{Stefenelli} submitted part of his paper to the Mathematical and Physical Class of the Munich Academy, apparently even expecting to see it included in the Academy's memoirs. Of course, this did not happen. We are not told whether the Academy gave a judgment on it. We would have been less severe in our review, or preferred to pass it over altogether, since as far as we know it seems to be the author's first attempt, were it not that his ill-considered attack on \textsc{Soldner}'s paper deserved a serious rebuke. For the author's own sake, it is regrettable that he was not led to give his attempt more time to mature before presenting it in public.\par\endgroup
